% ===== Appendix: corrections to the book, what we added, sources (after Questions; not presented) =====
\appendix
\partaccent{lavender}\gdef\idiottopicemoji{books}\gdef\idiottopicname{Հավելված}
\newcolumntype{L}[1]{>{\raggedright\arraybackslash}p{#1}}

\begin{frame}{Հավելված Ա. Ուղղումներ գրքի նկատմամբ (1/2)}
  \footnotesize
  \setlength{\tabcolsep}{4pt}
  \renewcommand{\arraystretch}{1.35}
  \begin{tabular}{@{}L{2.4cm}L{3.5cm}L{4.45cm}L{3.2cm}@{}}
    \toprule
    \textbf{Թեմա} & \textbf{Գրքում (2016)} & \textbf{Այսօր} & \textbf{Աղբյուր}\\
    \midrule
    Մեդուզաների քունը & «Մեդուզաները չեն քնում» & Cassiopea մեդուզան քնի նման վիճակ ունի, առանց ուղեղի & Nath et al., Curr. Biol., 2017\\
    Նորածինների REM & Քնի 80\%-ը & Մոտ 50\% (մեծահասակներ՝ մոտ 20\%) & Roffwarg et al., Science, 1966\\
    N1 փուլ & Ալֆա ալիքներ & Թետա ալիքներ. ալֆան հանգիստ արթնությանն է բնորոշ & AASM քնի գնահատման կանոններ\\
    \bottomrule
  \end{tabular}
\end{frame}

\begin{frame}{Հավելված Ա. Ուղղումներ գրքի նկատմամբ (2/2)}
  \footnotesize
  \setlength{\tabcolsep}{4pt}
  \renewcommand{\arraystretch}{1.35}
  \begin{tabular}{@{}L{2.4cm}L{3.5cm}L{4.45cm}L{3.2cm}@{}}
    \toprule
    \textbf{Թեմա} & \textbf{Գրքում (2016)} & \textbf{Այսօր} & \textbf{Աղբյուր}\\
    \midrule
    Մելատոնին & «Աճում է օրվա ընթացքում» & Ցերեկը ցածր է, բարձրանում է երեկոյան ու գիշերը & Դասագրքային\\
    «Սողունի ուղեղ» & Հարմար կրճատում & Եռամաս ուղեղի մոդելը հնացած է. մասերը միասին են էվոլյուցիա անցել & Cesario et al., Curr. Dir. Psychol. Sci., 2020\\
    Տեսողություն & «Ուղեղի մինչև 65\%-ը» & Գնահատականները շատ տարբեր են, սլայդում թիվ չենք դրել & \\
    \bottomrule
  \end{tabular}
\end{frame}

\begin{frame}{Հավելված Բ. Ինչ ենք ավելացրել գրքից դուրս}
  \begin{columns}[T]
    \begin{column}{0.48\textwidth}
      \begin{card}[colback=mint!30!cardwhite, fontupper=\scriptsize]
        {\blackfont Թվեր ու հետազոտություններ}\par\smallskip
        \begin{itemize}\setlength{\itemsep}{0.5mm}
          \item 86 միլիարդ նեյրոն՝ Azevedo et al., 2009
          \item Քնի կաթվածի հաճախականությունը՝ Sharpless \& Barber, 2011
          \item Քնի նորմաներն ըստ տարիքի՝ Hirshkowitz et al., 2015
        \end{itemize}
      \end{card}
    \end{column}
    \begin{column}{0.48\textwidth}
      \begin{card}[colback=peach!35!cardwhite, fontupper=\scriptsize]
        {\blackfont Մերն է, ոչ թե գրքից}\par\smallskip
        \begin{itemize}\setlength{\itemsep}{0.5mm}
          \item Ծովախտի դեմ խորհուրդները (բխում են զգայական կոնֆլիկտի վարկածից)
          \item Մեմերը. բովանդակությունը գրքից է
        \end{itemize}
      \end{card}
    \end{column}
  \end{columns}
\end{frame}

\begin{frame}{Հավելված Գ. Աղբյուրներ. գլուխ 1}
  \scriptsize
  Գրքի ծանոթագրություններից, որոնց վրա հիմնված են սլայդները\par\smallskip
  \begin{columns}[T]
    \begin{column}{0.48\textwidth}
      \begin{itemize}\setlength{\itemsep}{0.4mm}
        \item Dietz V. Spinal cord pattern generators for locomotion. Clin. Neurophysiol., 2003
        \item Ebenholtz S. M. et al. The possible role of nystagmus in motion sickness. Aviat. Space Environ. Med., 1994
        \item Weiler R. M. Olfaction and taste. J. Health Educ., 1999
        \item Wrangham R. Catching Fire: How Cooking Made Us Human. 2009
        \item Rechtschaffen A. et al. Physiological correlates of prolonged sleep deprivation in rats. Science, 1983
        \item Tononi G., Cirelli C. Perchance to prune. Sci. Am., 2013
      \end{itemize}
    \end{column}
    \begin{column}{0.48\textwidth}
      \begin{itemize}\setlength{\itemsep}{0.4mm}
        \item Gujar N. et al. Sleep deprivation amplifies reactivity of brain reward networks. J. Neurosci., 2011
        \item Siegel J. M. Sleep viewed as a state of adaptive inactivity. Nat. Rev. Neurosci., 2009
        \item Daan S. et al. Ground squirrels sleep during arousals from hibernation. Neurosci. Lett., 1991
        \item Brooks P. L., Peever J. H. Mechanisms responsible for REM sleep paralysis. J. Neurosci., 2012
        \item Driver H. S., Shapiro C. M. Parasomnias. BMJ, 1993
        \item Davis M. The role of the amygdala in fear and anxiety. Annu. Rev. Neurosci., 1992
      \end{itemize}
    \end{column}
  \end{columns}
\end{frame}
